\chapter{装进 VPS}
\chapterlead{没有物理接触的安装，全部通过远程控制台完成。本章在租用的一台 VPS 上部署 NixOS，说明镜像选择、远程救援的前提，以及为什么先要把回滚路径准备好。}

\section{起步：选一台能装自定义系统的 VPS}

三点条件必须满足：商家提供控制台（VNC 或串口），允许挂载自有 ISO，或允许从网络启动。缺少其中任何一项，都无法安装非预装系统。

\begin{booktable}{\textwidth}{J{0.846} J{1.154}}
\toprule
能力 & 用途 \\
\midrule
控制台访问 & 看不到屏幕时仍能操作 \\
挂载 ISO & 启动安装介质 \\
救援模式 & 系统起不来时进入临时环境 \\
快照 & 改动前的退路 \\
\bottomrule
\end{booktable}

\section{用 kexec 从现有系统切换}

许多商家不提供 ISO 挂载，但允许在现有 Linux 上执行命令。此时可以用 kexec 直接把 NixOS 安装环境加载进内存，无需重启到固件。

\begin{shellsession}
$ curl -L -o nixos.iso https://releases.nixos.org/nixos/26.11/nixos-26.11-x86_64-linux.iso
$ curl -L -o kexec.tar.xz https://github.com/nix-community/nixos-images/releases/latest/download/nixos-kexec-installer-x86_64-linux.tar.gz
$ tar -xf kexec.tar.xz
$ sudo ./kexec_nixos
\end{shellsession}

执行后机器会切换到安装环境，SSH 端口保持开放，用同一地址重新登录即可。这条路依赖商家的内核允许 kexec，绝大多数 KVM 机型允许。

\begin{warnbox}[先确认能回去]
kexec 之后，原系统不再运行。切换前记录原系统的网络配置（IP、网关、DNS）与 SSH 端口，并确认商家控制台可用。若安装失败且无法进控制台，只能请商家重装，数据全部丢失。
\end{warnbox}

\section{远程分区与安装}

进入安装环境后，流程与第二章一致，只是磁盘名与网络配置不同。

\begin{shellsession}
$ lsblk
$ sudo parted /dev/vda -- mklabel gpt
$ sudo parted /dev/vda -- mkpart ESP fat32 1MiB 512MiB
$ sudo parted /dev/vda -- set 1 esp on
$ sudo parted /dev/vda -- mkpart root ext4 512MiB 100%
$ sudo mkfs.fat -F 32 -n boot /dev/vda1
$ sudo mkfs.ext4 -L nixos /dev/vda2
$ sudo mount /dev/disk/by-label/nixos /mnt
$ sudo mkdir -p /mnt/boot
$ sudo mount /dev/disk/by-label/boot /mnt/boot
$ sudo nixos-generate-config --root /mnt
\end{shellsession}

服务器配置的差别在于网络与访问方式。云主机的网络通常由 DHCP 配置，保留默认即可；静态 IP 需要写进配置。

\begin{nixcode}
networking = {
  hostName = "srv1";
  useDHCP = false;
  interfaces.ens3.useDHCP = false;
  interfaces.ens3.ipv4.addresses = [{
    address = "203.0.113.10";
    prefixLength = 24;
  }];
  defaultGateway = "203.0.113.1";
  nameservers = [ "1.1.1.1" "8.8.8.8" ];
};

services.openssh = {
  enable = true;
  settings.PasswordAuthentication = false;
};
users.users.weiwen.openssh.authorizedKeys.keys = [
  "ssh-ed25519 AAAA... weiwen@laptop"
];
\end{nixcode}

把公钥写进配置是关键一步：安装完成后可以直接免密登录。密码登录关掉，减少被暴力猜测的暴露面。

\section{第一次远程登录}

\begin{shellsession}
$ ssh weiwen@203.0.113.10
$ nixos-version
$ sudo nixos-rebuild switch
\end{shellsession}

三条命令分别确认能登录、版本正确、重建流程可用。第三条尤其重要：它是后续所有改动的通道，必须先在无风险的状态下验证一遍。

\begin{notebox}[改 SSH 配置的顺序]
修改 SSH 相关配置时，保持一个已登录的会话不要断开，另开一个终端测试新连接。确认新连接可用后再关闭旧会话。这样即使配置写错，也还有一条活着的连接可以改回去。
\end{notebox}

\section{把配置放进版本控制}

服务器配置应当从第一天就在 Git 里。

\begin{shellsession}
$ sudo cp /etc/nixos/hardware-configuration.nix ~/srv1/
$ git init ~/srv1 && cd ~/srv1
$ git add . && git commit -m "initial server config"
$ git remote add origin git@example.com:weiwen/srv1.git
$ git push -u origin main
\end{shellsession}

此后 \texttt{/etc/nixos} 可以指向仓库中的文件，或反过来把仓库检出到家目录、用 \texttt{-{}-flake} 指向它。前者改动即时可见，后者历史清晰，两种都常见。

\section{练习}

\begin{enumerate}
  \item 列出所租 VPS 支持的三项能力：控制台、ISO 挂载、快照。
  \item 用 kexec 切换到安装环境，记录切换前后 IP 与 SSH 端口是否变化。
  \item 用静态 IP 完成一次安装，说明 \texttt{useDHCP} 与 \texttt{interfaces} 的关系。
  \item 关闭密码登录后，用两个终端验证“保留旧会话、测试新连接”的流程。
  \item 把服务器配置推到远端仓库，在另一台机器克隆并检查能否读到全部文件。
\end{enumerate}

\begin{checkpoint}
本章结束时，应能：
\begin{itemize}
  \item 判断一台 VPS 能否安装自定义系统；
  \item 用 kexec 或 ISO 挂载进入安装环境，知道失败时的退路；
  \item 写出服务器所需的网络与 SSH 配置；
  \item 按“保留旧会话”的顺序修改远程访问配置；
  \item 把服务器配置纳入版本控制。
\end{itemize}
\end{checkpoint}
